\begin{env}[5.1.6]
\label{IV.2.5.1.6-fr}
Vamos a examinar ahora las propiedades más especiales ligadas a hipótesis
noetherianas.

Recordemos $(\mathbf{0},16.2.3)$ que un anillo \emph{local noetheriano} $A\ne0$ tiene
dimensión \emph{finita}, igual al número mínimo de generadores de un ideal de
definición de $A$; para todo ideal primo $\mathfrak p$ de $A$, la altura de
$\mathfrak p$, igual a $\dim(A_{\mathfrak p})$, es por tanto también finita.
Estas propiedades, junto con (5.1.2), (5.1.3) y (5.1.4), muestran que:
\end{env}

\begin{proposition}[5.1.7]
\label{IV.2.5.1.7-fr}
Para toda parte cerrada no vacía $Y$ de un preesquema localmente noetheriano
$X$, $\operatorname{codim}(Y,X)$ es finita. Si $X$ es noetheriano y afín e
$Y$ una parte cerrada irreducible de $X$, $\operatorname{codim}(Y,X)$ es
igual al número mínimo de secciones $s_i$ de $\mathcal O_X$ sobre $X$ tales
que $Y$ sea una componente irreducible del conjunto de los $x\in X$ tales que
$s_i(x)=0$ para todo $i$.
\end{proposition}

\begin{corollary}[5.1.8]
\label{IV.2.5.1.8-fr}
Sean $X$ un preesquema localmente noetheriano, $\mathcal L$ un
$\mathcal O_X$-módulo invertible y $f$ una sección de $\mathcal L$ sobre
$X$. Entonces toda componente irreducible del conjunto $Z$

\oldpage[IV-2]{88}

de los $x\in X$ tales que $f(x)=0$ es de codimensión $\leq1$ en $X$; es de
codimensión $1$ si $Z$ no contiene ninguna componente irreducible de $X$.
\end{corollary}

\begin{proof}
Es posible limitarse al caso en que $\mathcal L=\mathcal O_X$. Si $y$ es un
punto genérico de una componente irreducible $Y$ de $Z$, el ideal $(f_y)$ de
$\mathcal O_{X,y}$ debe ser tal que $\mathcal O_{X,y}/(f_y)$ no tenga más que
un ideal primo, lo cual significa que $f_y$ genera un ideal de definición del
anillo local noetheriano $\mathcal O_{X,y}$; se tiene, pues,
$\operatorname{codim}(Y,X)\leq1$ (5.1.7); si $Z$ no contiene ninguna
componente irreducible de $X$, no se puede tener
$\operatorname{codim}(Y,X)=0$ en virtud de $(\mathbf{0},14.2.1)$.
\end{proof}

\begin{proposition}[5.1.9]
\label{IV.2.5.1.9-fr}
Sean $X$ un preesquema localmente noetheriano y $Y$ una parte cerrada de $X$.
Si $x\in Y$ es tal que $\mathcal O_{X,x}$ es un anillo catenario, se tiene
\begin{align}
\operatorname{codim}_x(Y,X)
&=\dim(\mathcal O_{X,x})\setmin\operatorname{codim}(\overline{\{x\}},Y)\notag\\
&=\dim(\mathcal O_{X,x})\setmin\dim(\mathcal O_{Y,x})
\tag{5.1.9.1}\label{IV.2.5.1.9.1-fr}
\end{align}
\end{proposition}

\begin{proof}
Sean $Y_i$ ($1\leq i\leq n$) las componentes irreducibles de $Y$ que
contienen a $x$ (que son en número \emph{finito}, puesto que $Y$ es localmente
noetheriano), y sea $y_i$ el punto genérico de $Y_i$. Si se pone
$A=\mathcal O_{X,x}$, y si $\mathfrak p_i$ es el ideal primo de $A$
correspondiente a $Y_i\cap\operatorname{Spec}(A)$, las partes cerradas
irreducibles de $Y$ que contienen a $x$ se corresponden biunívocamente con
los ideales primos de $A$ que contienen uno de los $\mathfrak p_i$; por tanto,
$\dim(\mathcal O_{Y,x})=\sup_i\dim(A/\mathfrak p_i)$; por otra parte, se tiene
$\mathcal O_{X,y_i}=A_{\mathfrak p_i}$, y la hipótesis sobre $A$ implica
$\dim(A)=\dim(A/\mathfrak p_i)+\dim(A_{\mathfrak p_i})$ $(\mathbf{0},16.1.4)$; la
conclusión resulta entonces de la relación
$\operatorname{codim}_x(Y,X)=\inf_i(\dim(\mathcal O_{X,y_i}))$ ((5.1.2) y
$(\mathbf{0},14.2.6)$).
\end{proof}
